\begin{claim}{D.27}{Proved}{The full-profinite rank divisibility and prime-order lifting test}
For a torsion-free profinite extension
\[
 1\longrightarrow\widehat F_r\longrightarrow E
   \longrightarrow G\longrightarrow1,
 \qquad 2\le r<\infty,\quad |G|<\infty,
\]
the integer $|G|$ divides $r-1$. In particular a torsion-free
finite extension of $\widehat F_2$ equals its prescribed kernel.
Every homomorphism $C_p\to\Out(\widehat F_2)$, for any prime $p$,
has a homomorphic lift to $\Aut(\widehat F_2)$; a faithful one has
a lift of order exactly $p$. This last assertion holds in the full
profinite category, without an invariant discretization and without
first passing to a pro-$p$ quotient.
\end{claim}
\begin{proof}
We give the normal-subgroup version of Jarden's index argument
\cite[Theorems 1.2 and 4.1, pp.~496--498]{R-Jarden1982}, with the
quotient comparisons written out. Put $F=\widehat F_r$. For every
prime $q$, $\operatorname{cd}_q(F)=1$. The torsion-free finite-index
dimension theorem of Serre \cite[\S1]{R-SerreCD}, whose torsion
hypothesis applies to $E$ itself, gives $\operatorname{cd}_q(E)=1$.
The same is true of every open subgroup of $E$. Thus these groups
are projective profinite groups. Here we use the precise
projectivity criterion recalled in \cite[\S2]{R-Jarden1982}:
a profinite group is projective exactly when its cohomological
dimension at every prime is at most one; equivalently, all its
finite embedding problems have weak solutions. This is the
criterion for the full profinite category, including finite
embedding problems whose kernels are nonabelian.

Fix a prime divisor $p$ of $|G|$, and let $p^a$ be its exact
$p$-part. Choose a Sylow $p$-subgroup $G_p\le G$ and write
$E_p$ for its inverse image in $E$. Define $O^p(J)$, for a
profinite group $J$, as the kernel of its maximal pro-$p$ quotient.
The characteristic subgroup $O^p(F)$ of $F$ is normal in $E_p$.
There are exact continuous quotient maps
\[
 1\longrightarrow F/O^p(F)\longrightarrow E_p/O^p(F)
 \longrightarrow G_p\longrightarrow1.
\]
The middle group is pro-$p$, since both its kernel and quotient
are pro-$p$. Therefore $O^p(E_p)\le O^p(F)$. Conversely the
restriction to $F$ of the maximal pro-$p$ quotient of $E_p$ is a
pro-$p$ quotient of $F$, so $O^p(F)\le O^p(E_p)$. This proves
\[
 O^p(E_p)=O^p(F),\qquad
 [E_p^{(p)}:F^{(p)}]=p^a.                         \tag{D.27a}
\]

The quotient $E_p^{(p)}$ is projective as a pro-$p$ group. Indeed a
finite embedding problem with $p$-group target and covering group
pulls back to one for $E_p$; its projective solution factors through
$E_p^{(p)}$ because its target is a $p$-group. A projective pro-$p$
group is free pro-$p$. It has finite rank $d_p$ because $E_p$ is
finitely generated. The subgroup $F^{(p)}$ in (D.27a) is free
pro-$p$ of rank $r$, with its full induced topology, as the exact
quotient construction proves. Finite-coset Schreier rewriting
therefore gives
\[
                         r-1=p^a(d_p-1).          \tag{D.27b}
\]
Consequently each exact prime-power divisor $p^a$ of $|G|$ divides
$r-1$. These divisors are pairwise coprime, proving the first
assertion. At $r=2$, this gives $|G|=1$.

Now let $\rho:C_p\to\Out(\widehat F_2)$ be a prescribed continuous
outer action, with the finite group given its discrete topology.
Form the pullback
\[
 E_\rho=\{(a,g)\in\Aut(\widehat F_2)\times C_p:
                         [a]=\rho(g)\}.
\]
This is a profinite group. Centre-freeness identifies its kernel
over $C_p$ with $\Inn(\widehat F_2)=\widehat F_2$ by the injective
map $x\mapsto(\Inn(x),1)$. If $E_\rho$ were torsion-free, the
already proved rank-two assertion would force its quotient to be
trivial, contradicting $|C_p|=p$. Hence $E_\rho$ contains a
nonidentity torsion element $v$. Its finite cyclic subgroup
intersects the torsion-free kernel trivially, so projection embeds
it in $C_p$. Its order is therefore exactly $p$, and it projects
isomorphically onto $C_p$. The inverse of that projection is a
section with the prescribed labels. Composing with the first
projection $E_\rho\to\Aut(\widehat F_2)$ gives the asserted lift.
For faithful $\rho$, the lift is faithful because its outer
projection is. No choice of conjugacy class of this section is
claimed; the argument proves existence.

Applied to a prime-order element of $GT_0$, this gives a genuine
finite-order loop-preserving automorphism representing its outer
class. It does not impose an unproved restriction on its multiplier,
or identify its fixed subgroup. In particular it neither solves
the residual norm equation D.25 nor proves conjugacy to the
geometric involution. It does prove that the nontrivial
torsion-free recognition test of D.22 first reaches free rank
three, where (D.27b) permits index two and D.26 supplies an
explicit unmarked example.
\end{proof}
